% ═══════════════════════════════════════════════════════════════
% QR-BLATT - Minitest s(t)-Diagramme | Physik Klasse 6
% Drucken: 1x pro Tisch oder an Tafel zeigen
% ═══════════════════════════════════════════════════════════════

\documentclass[11pt, a4paper]{article}

\usepackage[utf8]{inputenc}
\usepackage[T1]{fontenc}
\usepackage[ngerman]{babel}
\usepackage{helvet}
\renewcommand{\familydefault}{\sfdefault}

\usepackage[margin=2cm]{geometry}
\usepackage{xcolor}
\usepackage{tcolorbox}
\usepackage{enumitem}
\usepackage{qrcode}
\usepackage[hidelinks]{hyperref}

% URL-Definition
\newcommand{\mturl}{https://devbox42.github.io/sciencesim/physik/kl06/02-bewegung/std-03-st-diagramme/MT-03-st-diagramme.html}

% Farben
\definecolor{physik}{HTML}{2c5aa0}
\definecolor{physik-hell}{HTML}{e8f0fa}

\tcbuselibrary{skins}

\setlength{\parindent}{0pt}
\setlength{\parskip}{8pt}

\pagestyle{empty}

\begin{document}

% ═══════════════════════════════════════════════════════════════
% TITEL
% ═══════════════════════════════════════════════════════════════

\begin{center}
{\Huge\bfseries\textcolor{physik}{Minitest: Weg-Zeit-Diagramme}}\\[0.5em]
{\large Physik Klasse 6 | 25 BE | ca. 20 Minuten}
\end{center}

\vspace{1.5em}

% ═══════════════════════════════════════════════════════════════
% QR-CODE
% ═══════════════════════════════════════════════════════════════

\begin{center}
\begin{tcolorbox}[
    colback=white,
    colframe=physik,
    boxrule=2pt,
    width=10cm,
    arc=0pt,
    halign=center
]
\vspace{0.5em}
% QR-Code (klickbar)
\href{\mturl}{\qrcode[height=6cm]{\mturl}}

\vspace{0.5em}
{\small\textcolor{gray}{Scanne diesen Code mit deinem Tablet}}
\end{tcolorbox}
\end{center}

\vspace{1.5em}

% ═══════════════════════════════════════════════════════════════
% ANLEITUNG
% ═══════════════════════════════════════════════════════════════

\begin{tcolorbox}[
    colback=physik-hell,
    colframe=physik,
    boxrule=1.5pt,
    arc=0pt,
    title={\large\bfseries So geht's:},
    coltitle=white,
    fonttitle=\sffamily
]

\begin{enumerate}[leftmargin=2em, itemsep=10pt, label={\textcolor{physik}{\Large\textbf{\arabic*.}}}]
    \item \textbf{Scanne den QR-Code} mit deinem Tablet.

    \item \textbf{Bearbeite alle 8 Aufgaben} sorgfältig.

    \item Drücke am Ende auf \textbf{„Auswerten"}.

    \item Dein Ergebnis wird \textbf{sofort} angezeigt!
\end{enumerate}

\end{tcolorbox}

\vspace{1.5em}

% ═══════════════════════════════════════════════════════════════
% HINWEIS
% ═══════════════════════════════════════════════════════════════

\begin{tcolorbox}[
    colback=white,
    colframe=physik,
    boxrule=1pt,
    arc=0pt
]
\textbf{Tipp:} Du darfst dein \textbf{Leseblatt (LB-03)} und das \textbf{Übungsblatt (AB-03)} verwenden!
\end{tcolorbox}

\vfill

% ═══════════════════════════════════════════════════════════════
% URL
% ═══════════════════════════════════════════════════════════════

\begin{center}
\textcolor{gray}{\small Falls der QR-Code nicht funktioniert:}\\[0.3em]
\texttt{\small devbox42.github.io/sciencesim/physik/kl06/02-bewegung/std-03-st-diagramme/MT-03-st-diagramme.html}
\end{center}

\end{document}
